\documentclass{ctexbook}
\usepackage{amsmath}
\title{反常积分的收敛判别法}
\begin{document}
	\section*{反常积分的收敛判别法}
	\subsection*{反常积分的Cauchy收敛原理}
	（Cauchy收敛原理）反常积分$\displaystyle\int_{a}^{+\infty}f(x)\,\mathrm{d} x$收敛的充分必要条件是，对$\forall \epsilon >0$,$\exists A_0 \ge a$，使对$\forall A,A' \ge A_0$，有
	
	\centerline{$|\displaystyle\int_{A}^{A'}f(x)\,\mathrm{d} x|<\epsilon$.}
	
	若$f(x)$在任意有限区间$[a,A]\subset [a,+\infty)$可积，且$\displaystyle\int_{a}^{+\infty}|f(x)|\,\mathrm{d} x$收敛，则称$\displaystyle\int_{a}^{+\infty}f(x)\,\mathrm{d} x$绝对收敛，$f(x)$在$[a,+\infty)$上绝对可积.
	
	若$\displaystyle\int_{a}^{+\infty}f(x)\,\mathrm{d} x$收敛但不绝对收敛，则称$\displaystyle\int_{a}^{+\infty}f(x)\,\mathrm{d} x$条件收敛，$f(x)$在$[a,+\infty)$上条件可积.
	
	由Cauchy收敛原理，可以得到，若$\displaystyle\int_{a}^{+\infty}f(x)\,\mathrm{d} x$绝对收敛，则其一定收敛.
	
	\subsection*{非负反常积分审敛法}
	（比较判别法）若在$[a,+\infty)$上恒有$0\le f(x)\le K\varphi (x)$，其中常数$K$满足$0<K<+\infty$，则有
	
	（1）当$\displaystyle\int_{a}^{+\infty}\varphi(x)\,\mathrm{d} x$收敛时，$\displaystyle\int_{a}^{+\infty}f(x)\,\mathrm{d} x$也收敛；
	
	（2）当$\displaystyle\int_{a}^{+\infty}f(x)\,\mathrm{d} x$发散时，$\displaystyle\int_{a}^{+\infty}\varphi(x)\,\mathrm{d} x$也发散.
	
	\hspace*{\fill}
	
	（比较判别法的极限形式）若在$[a,+\infty)$上恒有$f(x)\ge 0$和$\varphi (x)\ge 0$，且有$\lim\limits_{x\rightarrow +\infty}\dfrac{f(x)}{\varphi (x)}=l$，则有
	
	（1）$0<x<+\infty$时，$\displaystyle\int_{a}^{+\infty}f(x)\,\mathrm{d} x$与$\displaystyle\int_{a}^{+\infty}\varphi(x)\,\mathrm{d} x$具有相同的敛散性；
	
	（2）$x=0$时，$\displaystyle\int_{a}^{+\infty}\varphi(x)\,\mathrm{d} x$收敛时，$\displaystyle\int_{a}^{+\infty}f(x)\,\mathrm{d} x$也收敛；
	
	（3）$x=+\infty$时，$\displaystyle\int_{a}^{+\infty}\varphi(x)\,\mathrm{d} x$发散时，$\displaystyle\int_{a}^{+\infty}f(x)\,\mathrm{d} x$也发散.
	
	\hspace*{\fill}
	
	（Cauchy判别法）设在$[a,+\infty)\subset[0,+\infty)$上，有$f(x)\ge 0$，$0<K<+\infty$，则有
	
	（1）若$f(x)\le \dfrac{K}{x^p}$，且$p>1$，则$\displaystyle\int_{a}^{+\infty}f(x)\,\mathrm{d} x$收敛；
	
	（2）若$f(x)\ge \dfrac{K}{x^p}$，且$p\le1$，则$\displaystyle\int_{a}^{+\infty}f(x)\,\mathrm{d} x$发散.
	
	\hspace*{\fill}
	
	（Cauchy判别法的极限形式）设在$[a,+\infty)\subset[0,+\infty)$上，有$f(x)\ge 0$，且
	
	\centerline{$\lim\limits_{x\rightarrow +\infty}x^p f(x)=l$}
	
	（1）若$0\le l<+\infty$，且$p>1$，则$\displaystyle\int_{a}^{+\infty}f(x)\,\mathrm{d} x$收敛；
	
	（2）若$0< l\le +\infty$，且$p\le1$，则$\displaystyle\int_{a}^{+\infty}f(x)\,\mathrm{d} x$发散.
	
	\subsection*{一般反常积分审敛法}
	（A-D判别法）若下列任一条件满足，则$\displaystyle\int_{a}^{+\infty}f(x)g(x)\,\mathrm{d} x$收敛：
	
	（1）（Abel判别法）$\displaystyle\int_{a}^{+\infty}f(x)\,\mathrm{d} x$收敛，$g(x)$在$[a,+\infty)$上单调有界；
	
	（2）（Dirichlet判别法）$F(A)=\displaystyle\int_{a}^{A}f(x)\,\mathrm{d} x$在$[a,+\infty)$上有界，$g(x)$在$[a,+\infty)$上单调且$\lim\limits_{x\rightarrow +\infty}g(x)=0$.
	
	\subsection*{无界反常积分审敛法}
	令下面讨论的无界反常积分的函数$f(x)$在$[a,b]$上只有$x=b$一个奇点.
	
	（Cauchy收敛原理）反常积分$\displaystyle\int_{a}^{b}f(x)\,\mathrm{d} x$收敛的充分必要条件是，对$\forall \epsilon >0$,$\exists \delta > 0$，使对$\forall \eta,\eta' \in (0,\delta)$，有
	
	\centerline{$|\displaystyle\int_{b-\eta}^{b-\eta'}f(x)\,\mathrm{d} x|<\epsilon$.}
	
	\hspace*{\fill}
	
	（Cauchy判别法）设在$[a,b)$上，恒有$f(x)\ge0$，而当$x$在$b$的某个左邻域上$[b-\eta_0,b)$时，存在$0<K<+\infty$，使得
	
	（1）$f(x)\le \dfrac{K}{(b-x)^p}$，且$p<1$，则$\displaystyle\int_{a}^{b}f(x)\,\mathrm{d} x$收敛；
	
	（2）$f(x)\ge \dfrac{K}{(b-x)^p}$，且$p\ge1$，则$\displaystyle\int_{a}^{b}f(x)\,\mathrm{d} x$发散.
	
	\hspace*{\fill}
	
	（Cauchy判别法的极限形式）设在$[a,b)$上，恒有$f(x)\ge0$，且
	
	\centerline{$\lim\limits_{x\rightarrow b-}(b-x)^p f(x)=l$}
	
	（1）若$0\le l<+\infty$，且$p<1$，则$\displaystyle\int_{a}^{b}f(x)\,\mathrm{d} x$收敛；
	
	（2）若$0< l\le +\infty$，且$p\ge1$，则$\displaystyle\int_{a}^{b}f(x)\,\mathrm{d} x$发散.
	
	\hspace*{\fill}
	
	（A-D判别法）若下列任一条件满足，则$\displaystyle\int_{a}^{b}f(x)g(x)\,\mathrm{d} x$收敛：
	
	（1）（Abel判别法）$\displaystyle\int_{a}^{b}f(x)\,\mathrm{d} x$收敛，$g(x)$在$[a,b)$上单调有界；
	
	（2）（Dirichlet判别法）$F(\eta)=\displaystyle\int_{a}^{b-\eta}f(x)\,\mathrm{d} x$在$(0,b-a]$上有界，$g(x)$在$[a,b)$上单调且$\lim\limits_{x\rightarrow b-}g(x)=0$.
\end{document}